\section{Post-Iteration Processing}
\label{sec:misc-details}
%% ===================================================================

The per-pixel iteration kernels produce a buffer of raw escape-time
counts---integers that record how many iterations each pixel survived
before escaping.  These counts are not directly displayable: they must
be mapped to colours, and the resulting image must be antialiased to
suppress the severe spatial aliasing that arises at deep zooms (where
intricate fractal features routinely fall below pixel scale).  Without
careful post-processing, the visual output would be a noisy, jagged
approximation of the underlying mathematics.

This section covers the three main post-iteration stages:
antialiasing via GPU supersampling (\cref{subsec:antialiasing}), the
palette-based coloring system (\cref{subsec:coloring}), and
iteration-count statistics gathered via parallel GPU reduction.  The
subsequent section (\cref{sec:render-thread-pool}) describes the
threading architecture that orchestrates these stages alongside the UI.

% -----------------------------------------------------------------
\subsection{Antialiasing via supersampling}
\label{subsec:antialiasing}

\FractalShark{} supports GPU-side supersampled antialiasing (SSAA).
The iteration kernels render at a resolution that is a multiple of
the output size; a separate \emph{antialiasing kernel} then
down-samples the result to the final pixel grid.

The supersampling factor is a compile-time template parameter
\code{Antialiasing}, explicitly instantiated for $1\times$, $2\times$,
$3\times$, and $4\times$.  For a factor~$k$, each output pixel
corresponds to a $k\times k$ block in the high-resolution iteration
buffer \code{OutputIterMatrix}.  The kernel loops over the $k^2$
samples, maps each iteration count to a palette colour
(\cref{subsec:coloring}), accumulates the RGB channels, and divides
by $k^2$ to produce an averaged \code{Color16} (16-bit per channel
RGBA) output.

\paragraph{Stream separation.}
The antialiasing kernel is launched on a dedicated CUDA stream
(\code{m\_Stream1}), distinct from the compute stream used by the
iteration kernels.  This allows the display thread to run
antialiasing and read back the colour buffer while the compute thread
continues issuing iteration work on its own default stream
(see also the \code{--default-stream per-thread}
compilation flag in \code{CudaDefaults.props}).

\paragraph{Kernel template.}
\begin{lstlisting}
template<typename IterType,
         uint32_t Antialiasing,
         bool ScaledColor>
__global__ void antialiasing_kernel(
    const IterType *OutputIterMatrix,
    uint32_t Width, uint32_t Height,
    AntialiasedColors OutputColorMatrix,
    Palette Pals,
    int local_color_width,
    int local_color_height,
    IterType n_iterations);
\end{lstlisting}
The \code{ScaledColor} flag selects between raw and scaled palette
indexing.  Pixels whose iteration count equals \code{n\_iterations}
(the Mandelbrot-set interior) are left black.

\paragraph{Limitations.}
Brute-force SSAA supersamples every pixel uniformly, which is
wasteful: most of the image consists of smooth-gradient regions where
a single sample suffices, while visual aliasing concentrates along the
boundary of the Mandelbrot set.  An edge-aware or adaptive
supersampling strategy---sampling densely only near detected
edges---would achieve comparable image quality at a fraction of the
cost.  The current uniform approach was chosen for implementation
simplicity on the GPU (all threads perform identical work, avoiding
divergence), but replacing it with an adaptive scheme is a natural
future improvement.

% -----------------------------------------------------------------
\subsection{Coloring and palette system}
\label{subsec:coloring}

After the antialiasing kernel maps iteration counts to colours, the
mapping itself is controlled by the palette system implemented in
\code{FractalPalette}.

\paragraph{Palette types.}
Five palettes are available, enumerated by
\code{FractalPaletteType}:
\code{Basic} (placeholder),
\code{Default} (rainbow transitions R$\to$Y$\to$G$\to$C$\to$B$\to$M),
\code{Patriotic} (red/blue),
\code{Summer} (red/green/cyan/white), and
\code{Random} (procedurally generated).
Palette data is pre-generated on the CPU (using smooth colour
transitions via \code{PalTransition}) and uploaded to the GPU as an
array of \code{Color16} values.

\paragraph{Bit-depth levels.}
Each palette is generated at six bit-depth levels, corresponding to
$2^5$, $2^6$, $2^8$, $2^{12}$, $2^{16}$, and $2^{20}$ colours
(32 to roughly one million).  The active depth is selected by
\code{m\_PaletteDepthIndex} (default: index~2, i.e.\ 256 colours)
and can be cycled at runtime.

\paragraph{Iteration-to-colour mapping.}
On the GPU, the palette index for a given iteration count is computed
as
\[
  \mathrm{palIndex} \;=\;
  \bigl(\mathtt{numIters} \gg \mathtt{palette\_aux\_depth}\bigr)
    \bmod \mathtt{local\_palIters},
\]
where \code{palette\_aux\_depth} (range 0--16) controls the iteration
compression---higher values produce wider colour bands---and
\code{local\_palIters} is the palette size at the active bit depth.
The modulo wrapping causes the palette to cycle when the iteration
count exceeds the palette length.

\paragraph{Palette rotation.}
A rotation offset \code{m\_PaletteRotate} is added to the iteration
count before palette lookup, allowing the user to phase-shift the
colour mapping for visual exploration without re-rendering.

\paragraph{Alternative coloring approaches.}
The palette-indexed scheme above is simple and fast but has known
limitations: discrete iteration counts produce visible banding, and
the palette cycle length is fixed.  Several alternative techniques
exist in the fractal-rendering community:
\begin{itemize}
  \item \textbf{Smooth (continuous) iteration count.}  By incorporating
    the final orbit magnitude into a fractional escape count
    ($n_{\text{smooth}} = n - \log_2\!\log_2|z_n| + \text{const}$),
    banding is eliminated and colour transitions become continuous.
  \item \textbf{Distance estimation.}  The analytic derivative
    $\partial z_n/\partial c$ (already computed for perturbation
    purposes) can yield an estimate of the distance from~$c$ to the
    Mandelbrot boundary.  Mapping distance to colour emphasises fine
    boundary structure and enables resolution-independent rendering of
    thin filaments.
  \item \textbf{Histogram equalization.}  Collecting a histogram of
    iteration counts across the image and mapping colours via the
    cumulative distribution produces uniform colour utilisation
    regardless of the iteration distribution---useful at deep zooms
    where most counts cluster in a narrow range.
  \item \textbf{Orbit-trap coloring.}  Measuring the closest approach
    of the orbit to a geometric shape (circle, cross, point) during
    iteration provides a distance metric unrelated to escape time,
    yielding distinctive visual effects.
\end{itemize}
\FractalShark{} currently implements only the palette-indexed method.
Integrating smooth iteration counts is feasible with minimal additional
complexity and would
pair naturally with the existing palette infrastructure; distance
estimation would leverage derivative data that the perturbation
pipeline already maintains.

% -----------------------------------------------------------------
\subsection{Iteration statistics via parallel reduction}
\label{subsec:reduction}

After the antialiasing kernel completes, a single reduction pass
extracts global statistics from the iteration buffer.  The
\code{max\_kernel} template reads every element of
\code{OutputIterMatrix} and computes three values simultaneously:
\begin{itemize}
  \item the minimum iteration count,
  \item the maximum iteration count,
  \item the sum of all iteration counts.
\end{itemize}
These are stored in a \code{ReductionResults} structure
(three \code{uint64\_t} fields: \code{Min}, \code{Max}, \code{Sum}).

The kernel uses shared-memory accumulators within each thread block
and atomic operations (\code{atomicMax}, \code{atomicMin},
\code{atomicAdd}) to merge block-level results into a single global
output.  It is launched with a $16\times16$ block grid and runs on
the same stream as the antialiasing kernel, so no additional
synchronization is required.  After the kernel completes, the host
copies back the \code{ReductionResults} via \code{cudaMemcpy} in
\code{ExtractItersAndColors}.

\paragraph{Current status.}
The reduction results are computed every frame but are not currently
used to drive any rendering decision (e.g.\ adaptive palette scaling
or histogram equalization).  They exist primarily for diagnostic
interest and as infrastructure that could support future coloring
enhancements.


% ===========================================================================
\section{Render Thread Pool and Command Queue}
\label{sec:render-thread-pool}
% ===========================================================================

The UI thread must remain responsive while fractal rendering proceeds.
\FractalShark{} achieves this through a \emph{command-queue} architecture:
the UI thread never mutates \code{Fractal} state directly.  Instead it wraps
mutations in lambdas and enqueues them into a work queue consumed by a pool of
worker threads.

\subsection{Architecture}

\begin{figure}[!htbp]
\centering
\begin{tikzpicture}[
    >=Stealth,
    box/.style={draw, rounded corners, minimum height=0.7cm, align=center,
                font=\small},
    ui/.style={box, fill=green!10, minimum width=2.0cm},
    queue/.style={box, fill=yellow!12, minimum width=2.0cm},
    worker/.style={box, fill=blue!10, minimum width=2.0cm},
    gl/.style={box, fill=purple!10, minimum width=2.0cm},
]
  \node[ui] (ui) {UI thread};
  \node[queue, right=0.7cm of ui] (eq) {Work queue};
  \node[worker, right=0.7cm of eq] (w) {Workers\\{\scriptsize($\times 4$)}};
  \node[worker, right=0.7cm of w] (calc) {\code{CalcFractal}};
  \node[gl, right=0.7cm of calc] (frame) {ProduceFrame};
  \node[gl, below=0.5cm of frame] (glc) {GL consumer};

  \draw[->] (ui) -- node[above, font=\scriptsize] {enqueue} (eq);
  \draw[->] (eq) -- node[above, font=\scriptsize] {dequeue} (w);
  \draw[->] (w) -- node[above, font=\scriptsize] {render} (calc);
  \draw[->] (calc) -- node[above, font=\scriptsize] {} (frame);
  \draw[->] (frame) -- (glc);

  \draw[->, dashed, gray] (eq.south) -- ++(0,-0.5) node[below, font=\scriptsize, gray, align=center] {supersede\\stale items};
\end{tikzpicture}
\caption[Render thread pool data flow]{Render thread pool data flow.
  The UI thread enqueues command+render work items.  Worker threads
  dequeue, execute the command lambda under a mutex, call
  \code{CalcFractal} to render into a private
  \code{Iters\-Memory\-Container}, then produce a frame for the OpenGL
  consumer thread.  Stale items in the queue are superseded by newer
  enqueues.}
\label{fig:render-pipeline}
\end{figure}

\code{RenderThreadPool} owns four worker threads (one per GPU renderer) and a
dedicated OpenGL consumer thread.  The work queue is a \code{std::deque} of
\code{RenderWorkItem} structs, each containing:

\begin{itemize}
  \item An optional \code{Command} lambda (\code{std::function<void(Fractal\&)>})
        executed under \code{m\_CalcFractalMutex} before rendering.
  \item An optional \code{Ptz} (coordinates) for the render; if absent, the
        worker snapshots the live \code{m\_Ptz} after command execution.
  \item Snapshotted algorithm, iteration count, dimensions, and dirty flags.
  \item A \code{CompletionPromise} signalled when the job finishes, enabling
        callers to \code{Wait()} if needed.
\end{itemize}

Workers execute a simple loop: dequeue an item, acquire a renderer and
\code{ItersMemoryContainer} from their respective pools, lock
\code{m\_CalcFractalMutex}, execute any command lambda, snapshot state, call
\code{CalcFractal}, then produce frames for the GL consumer.

\subsection{Work Item Structure}
\label{subsec:work-item-struct}

Each \code{RenderWorkItem} encapsulates a complete, self-contained
description of a rendering job.  Workers read exclusively from the work
item's fields, never from live \code{Fractal} members (with the
exception of the \code{FractalPtr} back-pointer used for reference
orbit access).  The struct comprises several categories of fields:

\begin{itemize}
  \item \textbf{Identity.}  A monotonic \code{SequenceNumber}
        (assigned at enqueue time) uniquely identifies the job and
        determines the display order in the GL consumer.  The
        \code{EnqueueGeneration} counter enables stale-render
        detection (\cref{subsec:two-render-paths}).
  \item \textbf{Coordinates.}  A \code{PointZoomBBConverter Ptz}
        snapshot captures the view centre, bounding box, and zoom
        factor.  If a \code{Command} lambda is present, the
        coordinates are re-snapshotted after command execution so that
        any mutations (pan, zoom) are reflected.
  \item \textbf{Render configuration.}  \code{Algorithm},
        \code{IterType}, \code{NumIterations},
        \code{IterationPrecision}, \code{ScrnWidth},
        \code{ScrnHeight}, and \code{GpuAntialiasing} are
        snapshotted from the \code{Fractal} at enqueue time.
        Three dirty flags (\code{ChangedWindow},
        \code{ChangedScrn}, \code{ChangedIterations}) are always
        set to \code{true} in the snapshot so that
        \code{CalcFractal} does not skip computation due to stale
        clean flags from a prior render.
  \item \textbf{Command and control.}  An optional
        \code{Command} lambda is executed under the calc-fractal
        mutex before rendering.  \code{MutationOnly} items skip the
        GPU pipeline entirely.  \code{Supersedable} controls whether
        a newer enqueue may discard this item.
  \item \textbf{Completion signalling.}  A
        \code{std::shared\_ptr<std::promise<void>{}>}
        (\code{CompletionPromise}) is signalled when the job
        finishes, whether by successful rendering, superseding, or
        error.  Callers may wait via the returned
        \code{RenderJobHandle}.
\end{itemize}

\subsection{ItersMemoryContainer Pool}
\label{subsec:iters-pool}

Each worker requires a private \code{ItersMemoryContainer} to hold
iteration results without contending on \code{m\_CurIters}.  A pool
of pre-allocated containers resides in
\code{Fractal::m\_ItersMemoryStorage}, guarded by a mutex and
condition variable.

\paragraph{Acquisition.}
\code{AcquireItersMemory} blocks under
\code{m\_ItersMemoryStorageLock} until the pool is non-empty, then
moves the back element out of the vector.  Because the pool is
initially seeded with enough containers for all workers, blocking
occurs only transiently---when all workers are active and no
container has yet been returned.

\paragraph{Return and validation.}
\code{ReturnIterMemory} checks that the container's dimensions and
antialiasing factor still match the current screen configuration.
Stale containers (from before a resize or antialiasing change) are
silently discarded rather than returned, which prevents a subsequent
worker from inheriting a mis-sized buffer.  Valid containers are
pushed back onto the vector and the condition variable is notified.

\paragraph{Worker-to-shared swap.}
After a successful render, the worker swaps its private
\code{ItersMemoryContainer} into \code{m\_CurIters} under the
calc-fractal mutex (if dimensions match).  This ensures that
facilities requiring CPU-side iteration data
(\code{SaveCurrentFractal}, \code{FindInterestingLocation}) see the
most recently completed render without requiring a separate direct-path
call.  If dimensions do not match---e.g.\ a resize occurred between
enqueue and completion---the swap is skipped and the old
\code{m\_CurIters} is retained.

\subsection{Mutex Ordering and Deadlock Prevention}
\label{subsec:mutex-ordering}

The pool architecture introduces several mutexes with overlapping
lifetimes.  The following lock-ordering discipline is maintained to
prevent deadlock:
\begin{enumerate}
  \item \code{m\_WorkQueueMutex} --- protects the work queue and
        sequence counter.
  \item \code{m\_CalcFractalMutex} --- serialises
        \code{CalcFractal} calls and shared \code{Fractal} state
        access.
  \item \code{RendererPool::m\_Mutex} --- guards renderer
        checkout/return.
  \item \code{m\_ItersMemoryStorageLock} --- guards the
        \code{ItersMemoryContainer} pool.
  \item \code{m\_BufferPoolMutex} --- guards the \code{Color16}
        frame-buffer pool.
  \item \code{FrameCompletionQueue::m\_Mutex} --- guards the
        completed-frame queue.
  \item \code{m\_DrainMutex} --- guards the drain condition
        variable.
\end{enumerate}
Tombstone frames are pushed \emph{outside} the work-queue lock to
avoid a nested lock between \code{m\_WorkQueueMutex} and the
\code{FrameCompletionQueue} mutex.  The in-flight counter
\code{m\_InFlightCount} is incremented atomically inside
\code{DequeueWorkItem} while the work-queue lock is still held,
closing a TOCTOU window in which \code{Drain()} could observe an
empty queue with zero in-flight count between pop and the
increment.

\subsection{Compute Stream and Display Stream Interaction}
\label{subsec:stream-interaction}

Each \code{GPURenderer} maintains two CUDA streams: a
\emph{compute stream} on which iteration kernels are launched,
and a \emph{display stream} (\code{m\_Stream1}) used by the
antialiasing and colour-extraction kernels
(\cref{subsec:antialiasing}).  The two-stream design permits
overlapped execution:

\begin{itemize}
  \item \textbf{Progressive frames.}
        \code{ProduceFrame} with \code{isFinal\!=\!false}
        synchronises on the display stream only
        (\code{SyncDisplayStream}), allowing the compute stream to
        continue issuing iteration work.  The antialiasing kernel
        reads whatever partial results have been written to the
        iteration buffer so far, producing a visually useful but
        incomplete image.
  \item \textbf{Final frames.}
        \code{ProduceFrame} with \code{isFinal\!=\!true}
        synchronises on the compute stream
        (\code{SyncComputeStream}) to guarantee that all iteration
        kernels have completed before colour extraction.
  \item \textbf{Compute-done notification.}
        Workers register a per-worker mutex/CV pair via
        \code{SetComputeDoneNotification}.  A CUDA host callback
        enqueued on the compute stream (\code{EnqueueComputeDoneCallback})
        sets an atomic flag and notifies the CV when all compute work
        finishes.  \code{WaitForGpuAndProduceProgressiveFrames}
        polls this flag with a one-second timeout, emitting
        progressive frames at each interval until the callback fires.
\end{itemize}

\subsection{Enqueue APIs}

Three enqueue variants serve distinct purposes:

\begin{center}
\begin{tabularx}{\textwidth}{lccX}
\hline
\textbf{API} & \textbf{Mutates} & \textbf{Renders} & \textbf{Use case} \\
\hline
\code{EnqueueCommand} & yes & yes & UI zoom, pan, palette changes, recalc \\
\code{EnqueueMutation} & yes & no  & Settings (algorithm, precision, perturbation type) \\
\code{EnqueueRender}   & no  & yes & Re-render current state; AutoZoomer pipeline \\
\hline
\end{tabularx}
\end{center}

\code{Enqueue\-Mutation} skips sequence-number assignment, renderer
acquisition, \code{Calc\-Fractal}, and frame production.  Its fast
path only acquires the calc-fractal mutex, executes the lambda, and
signals the promise.

\subsection{Queue Superseding and Generation Counter}

Rapid user input (e.g.\ holding a zoom key) can enqueue many renders of which
only the latest matters.  Two mechanisms discard stale work:

\begin{enumerate}
  \item \textbf{Queue-level superseding.}  When a new render item is enqueued,
        all earlier \emph{supersedable} render items still in the queue are
        removed.  Tombstone frames are pushed for their sequence numbers so the
        GL consumer can advance.  Mutations and non-supersedable items
        (AutoZoomer animation frames) are preserved.

  \item \textbf{Generation counter.}  Each supersedable enqueue increments a
        monotonic counter and stamps the item.  Workers check the counter
        \emph{before} acquiring a renderer: if a newer item has been enqueued
        since dequeue, the stale item is skipped with a tombstone.  This catches
        items that were already dequeued but not yet processing.
\end{enumerate}

AutoZoomer's pipelined animation frames set \code{Supersedable = false} so they
are never discarded.

\subsection{Two Rendering Paths}
\label{subsec:two-render-paths}

\code{CalcFractal} is invoked from two distinct paths with different iteration
data flow:

\paragraph{Render pool path (normal UI rendering).}
A worker acquires a private \code{ItersMemoryContainer} from a pool.  GPU
results are extracted via \code{RenderCurrent} into a \code{RenderFrame} and
forwarded to the GL consumer thread for texture upload.  The shared
\code{m\_CurIters} is \emph{not} updated.

\paragraph{Direct path (\code{CalcFractal(true)}).}
Used by \code{CrummyTest}, \code{AutoZoomer} Default/Max heuristic, and save
operations.  The caller constructs
\code{CalcContext\{m\_Ptz,\;m\_CurIters\}} so GPU results are written directly
into \code{m\_CurIters}.  The \code{drawFractal = true} flag triggers a
GPU\(\to\)CPU iteration copy at the end of \code{CalcFractalTypedIter} via
\code{RenderCurrent}.

\textbf{Invariant.}  Any code that reads \code{m\_CurIters} (saving images,
analysing iteration data) must use the direct path and call
\code{RenderThreadPool::Drain()} first to ensure no pool workers are active.
The pool path writes to \code{workerIters}, which is returned to the pool---\code{m\_CurIters} never sees that data.

\paragraph{Data-flow trace.}
In the pool path, the data-flow sequence proceeds as follows:
\begin{enumerate}
  \item The worker thread calls \code{AcquireItersMemory} to obtain a
        private \code{ItersMemoryContainer} from the pool.
  \item \code{RunCalcFractal} passes this container through a
        \code{CalcContext} and invokes \code{CalcFractal}, which
        dispatches GPU iteration kernels.  GPU results remain in
        device memory until extraction.
  \item \code{WaitForGpuAndProduceProgressiveFrames} polls the
        compute-done callback; while waiting, it periodically
        calls \code{ProduceFrame(isFinal\!=\!false)} to emit
        progressive (partial) frames.
  \item On compute completion, \code{ProduceFrame(isFinal\!=\!true)}
        calls \code{RenderCurrent} to run the antialiasing and
        colour-mapping kernels, copies the resulting \code{Color16}
        buffer from the GPU into a pooled frame buffer
        (\code{AcquireFrameBuffer}), and pushes a
        \code{RenderFrame} to the \code{FrameCompletionQueue}.
  \item The GL consumer thread dequeues the frame in sequence
        order via \code{TryPopNextInOrder}, uploads the
        \code{Color16} data as a \code{GL\_RGBA16} texture, draws
        a textured quad, and returns the buffer to the pool via
        \code{ReleaseFrameBuffer}.
\end{enumerate}
In the direct path, step~1 is replaced by constructing a
\code{CalcContext} around \code{m\_CurIters}; steps~3--5 are
replaced by an in-line \code{RenderCurrent} call at the end of
\code{CalcFractalTypedIter} that copies iteration data back to
CPU memory (the \code{drawFractal = true} branch).

\paragraph{Choosing between paths.}
The pool path is the default for all interactive rendering.  Its
private \code{ItersMemoryContainer} avoids write contention on
\code{m\_CurIters} and allows the GL consumer to display frames
without blocking the next render.  The direct path should be used
only when the caller requires CPU-side access to the final
iteration buffer---specifically for image saving
(\code{SaveCurrentFractal}), regression testing
(\code{CrummyTest}), and the AutoZoomer's heuristic analysis
modes that inspect iteration data.  Any such call must be preceded
by \code{Drain()} to ensure that no pool worker is concurrently
modifying shared state.

\paragraph{Performance implications.}
The pool path enables pipeline parallelism: while the GL consumer
uploads frame~$n$, a worker can already be computing frame~$n\!+\!1$.
Progressive frames provide visual feedback during long GPU
computations at the cost of one additional \code{RenderCurrent}
call per progress interval (default: 1\,s).  Progressive frames
synchronise via the display stream rather than the compute stream,
so the iteration kernels continue executing concurrently
(\cref{subsec:antialiasing}).

The direct path is inherently serial: the caller blocks until
\code{CalcFractal} completes, synchronises the compute stream
twice (once for iteration kernels, once for \code{RenderCurrent}),
and only then returns control.  It should therefore be avoided in
latency-sensitive UI code paths.

\subsection{MPIR Allocator Lifecycle}
\label{subsec:mpir-allocator-lifecycle}

MPIR memory allocation is controlled by a process-global function-pointer
triple installed via \code{mp\_set\_memory\_functions}.  Two custom allocators
are used to reduce allocation overhead during reference-orbit computation:

\begin{itemize}
  \item \textbf{\code{MPIRBoundedAllocator}} --- four $1\,\mathrm{MB}$
        stack-like blocks per thread (thread-local storage).  Allocs bump
        downward; frees decrement a per-block counter.  Used for short-lived
        temporaries in \code{CalcCpuHDR} and single-threaded reference-orbit
        paths.
  \item \textbf{\code{MPIRBumpAllocator}} --- backed by
        \code{GrowableVector<uint8\_t>}, grows dynamically.  Freed memory is
        recycled via a small LIFO cache.  Used for \code{SaveForReuse} orbit
        data that persists after computation.
\end{itemize}

\textbf{Lifetime constraint.}  All \code{mpf\_t} / \code{HighPrecision} objects
allocated under a custom allocator must be destroyed \emph{before}
\code{ShutdownTls()}.  After shutdown, \code{TlsInitialized} is false and
\code{mpf\_clear} falls back to the default allocator, which cannot recognise
pointers from the custom allocator.  Enclosing \code{mpf\_t} variables in
\code{\{\,\}} scope blocks ensures they are destroyed while the allocator TLS
is still active.

\subsection{Shutdown}

On application exit, \code{RenderThreadPool::Shutdown()} sets the shutdown flag,
clears the work queue (signalling completion on discarded items), and joins all
worker threads.  Workers exit immediately on seeing the shutdown flag, even if
items remain in the queue.  In-progress \code{CalcFractal} calls check
\code{m\_ShutdownFlag} during the GPU-wait loop and abort early.

% -----------------------------------------------------------------
\subsection{Coordinate management}
\label{subsec:coordinates}

The \code{PointZoomBBConverter} class manages the bidirectional
mapping between screen pixel coordinates and points on the complex
plane.  All internal coordinate quantities---centre, bounding-box
corners, zoom factor, radius---are stored as \code{HighPrecision}
values, which wrap MPIR's \code{mpf\_t} arbitrary-precision
floating-point type.  This allows the coordinate system to maintain
meaningful precision at extreme zoom depths where IEEE~754 types are
long exhausted.

\paragraph{Bounding box from centre and zoom.}
Given a centre $(c_x, c_y)$ and zoom factor~$Z$, the visible region
is
\[
  x \in \bigl[c_x - 2/Z,\; c_x + 2/Z\bigr],\qquad
  y \in \bigl[c_y - 2/Z,\; c_y + 2/Z\bigr].
\]
The constant~2 (the \code{Factor} member) sets the initial
un-zoomed viewport to the range $[-2,2]$ in both axes.

\paragraph{Screen-to-complex and complex-to-screen transforms.}
Pixel coordinates are mapped linearly:
\[
  x_{\text{calc}} = \frac{(x_{\text{px}} - x_0)\,(x_{\max}-x_{\min})}
                         {W \cdot k},
\]
where $W$ is the screen width and $k$ is the antialiasing factor.
The $y$~axis is inverted (screen~$y$ increases downward; imaginary
axis increases upward).  The inverse mapping is used to convert
fractal coordinates back to pixel positions.

\paragraph{Per-pixel stepping for the GPU.}
For GPU kernel launches, the converter computes per-pixel deltas
$\Delta x = (x_{\max}-x_{\min})/(W \cdot k)$ and
$\Delta y = (y_{\max}-y_{\min})/(H \cdot k)$ at full
\code{HighPrecision}.  These are then down-converted to the
kernel's working type (e.g.\ \code{double}, \code{HDRFloat}) for
per-thread coordinate calculation on the GPU.

\paragraph{Down-conversion via \code{FillGpuCoords}.}
The \code{FillGpuCoords} template bridges the arbitrary-precision
coordinate system and the GPU kernel's working type.  It invokes
\code{GetDeltaX} and \code{GetDeltaY} to obtain the per-pixel deltas
at full \code{HighPrecision}, then calls \code{FillCoord} four times
to convert $(x_{\min},\; y_{\min},\; \Delta x,\; \Delta y)$ into the
kernel's type parameter~\code{T}.  Type-specific overloads of
\code{FillCoord} handle the conversion:
\begin{itemize}
  \item \textbf{\code{double} / \code{float}.}  A single
        \code{Convert<HighPrecision, T>} truncation.  Sufficient when the
        per-pixel delta exceeds the machine epsilon of the target type
        (roughly $2^{-52}$ for \code{double}, $2^{-23}$ for \code{float}).
  \item \textbf{\code{MattDbldbl} / \code{MattDblflt}.}  A two-component
        head--tail split: the head stores the leading-order truncation
        and the tail stores the residual $(\text{src} - \text{head})$,
        effectively doubling the mantissa width.
  \item \textbf{Quad types (\code{MattQDbldbl}, \code{MattQFltflt}).}
        Four successive residual subtractions extend the effective mantissa
        to roughly $4\times$ the width of the scalar type.
  \item \textbf{\code{HDRFloat<T>}.}  Direct cast from
        \code{HighPrecision} followed by \code{HdrReduce}, which
        normalises the mantissa--exponent representation so that the
        mantissa lies in $[0.5,1)$.
\end{itemize}
Because the conversion is applied identically to both the bounding-box
origin and the per-pixel delta, the GPU kernel can reconstruct any
pixel's complex-plane coordinate as
$c_{\text{px}} = x_{\min} + \text{px}\cdot\Delta x$ using only the
working type, with no further access to arbitrary-precision quantities.

\paragraph{The antialiasing factor in delta computation.}
\label{par:aa-delta}
The antialiasing factor~$k$ (\code{GpuAntialiasing}) enters the
per-pixel delta as a divisor in the denominator:
\[
  \Delta x = \frac{x_{\max} - x_{\min}}{W \cdot k}, \qquad
  \Delta y = \frac{y_{\max} - y_{\min}}{H \cdot k}.
\]
Multiplying the effective resolution by~$k$ shrinks the step size so
that $k^2$ sub-pixel samples are computed per output pixel.  Within
\code{PointZoomBBConverter}, the screen-to-complex conversion
functions \code{XFromScreenToCalc} and \code{YFromScreenToCalc}
accept \code{antialiasing} as a parameter and compute the origin
offset consistently:
\[
  \text{OriginX} = \frac{W \cdot k}{x_{\max} - x_{\min}}
                   \cdot (-x_{\min}),
\]
ensuring that screen coordinate $(0,0)$ maps to $(x_{\min}, y_{\max})$
regardless of the supersampling level.  The factor~$k$ therefore
propagates through the entire coordinate chain without special-case
logic: the iteration kernels operate on a
$(W\!\cdot\! k) \times (H\!\cdot\! k)$ grid, and the antialiasing
kernel (\cref{subsec:antialiasing}) down-samples the result.

\paragraph{Zoom operations.}
Zooming recomputes the bounding box around the current centre:
$x'_{\min} = c_x - r/d$, $x'_{\max} = c_x + r/d$ (and similarly
for~$y$), where $d$ is the zoom divisor.  All intermediate arithmetic
is performed at the working precision of the \code{HighPrecision}
members, which is dynamically adjusted via \code{SetPrecision} as
the zoom deepens.

\paragraph{Precision selection heuristics.}
\label{par:precision-heuristics}
The function \code{PrecisionCalculator::GetPrecision} determines how
many bits of mantissa the \code{HighPrecision} members require at a
given zoom depth.  It converts the bounding-box deltas $\Delta x$ and
$\Delta y$ to \code{HDRFloat<double>}, extracts the base-2 exponent
of each (via \code{getExp} or \code{std::frexp}, depending on the
numeric type), and takes the larger absolute exponent as the baseline.
Two additive padding constants control the final value:
\begin{itemize}
  \item \textbf{Standard precision.}
        $\text{bits} = |\text{exp}| + 120$
        (\code{AuthoritativeMinExtraPrecisionInBits}).
        The 120-bit padding provides roughly 36 decimal guard digits
        beyond the last significant bit of the per-pixel delta, which
        suffices for single-use reference orbit computation and
        coordinate arithmetic.
  \item \textbf{Reuse precision.}
        $\text{bits} = |\text{exp}| + 800$
        (\code{AuthoritativeReuseExtraPrecisionInBits}).
        When \code{RequiresReuse} is true (the orbit will be
        saved and reused across views), an additional 680 guard bits
        ensure that the reference orbit remains accurate even after
        recentring to a nearby coordinate.
\end{itemize}
\code{SetPrecision} propagates this value to every
\code{HighPrecision} member of the converter---$x_{\min}$,
$x_{\max}$, $y_{\min}$, $y_{\max}$, $c_x$, $c_y$, and
$Z$---via \code{precisionInBits}, and also updates MPIR's default
precision so that newly constructed \code{mpf\_t} temporaries inherit
the same width.  Because this call reallocates the underlying limb
arrays, it is performed once per zoom-level change rather than per
frame.

\paragraph{Maintaining \code{HighPrecision} throughout coordinate chains.}
\label{par:hp-invariant}
A deliberate design choice is that all coordinate quantities internal to
\code{PointZoomBBConverter} remain as \code{HighPrecision}
(\code{mpf\_t}) values throughout their lifecycle.  Zoom operations,
recentring, aspect-ratio correction (\code{SquareAspectRatio}), and
bounding-box recomputation all operate on the arbitrary-precision
members directly, using MPIR's in-place arithmetic (e.g.\
\code{subFrom}, \code{divFrom\_ui}, \code{mulFrom}).  Down-conversion
to a finite-precision type occurs only at the final stage when
\code{FillGpuCoords} prepares values for GPU upload (described
above under ``Down-conversion'').  This architecture ensures that
accumulated rounding errors from successive zoom and pan operations do
not degrade coordinate accuracy: each operation is computed at the full
working precision determined by \code{SetPrecision}, and only the
final GPU-facing values are truncated.

% -----------------------------------------------------------------
\subsection{GPU memory allocation with host fallback}
\label{subsec:gpu-mem-fallback}

Reference orbits at deep zoom levels can be extremely large---tens of
billions of iterations, each storing several floating-point
values---and may exceed the GPU's physical VRAM (typically 8--24\,GB
on consumer hardware).  However, the host system often has
substantially more RAM available.  Rather than aborting the render
when GPU memory is exhausted, \FractalShark{} uses a two-tier
fallback strategy.

\paragraph{Allocation cascade.}
Each large GPU allocation follows a try-device-then-host pattern:
\begin{enumerate}
\item Attempt \code{cudaMalloc} (device memory, full GPU bandwidth).
\item If that fails, attempt \code{cudaMallocHost}
      (pinned host memory, accessible from GPU via PCIe but at
      reduced bandwidth).
\item If both fail, mark the allocation as invalid and set
      its size to zero.
\end{enumerate}
A boolean flag (e.g.\ \code{AllocHost} in
\code{GPUPerturbSingleResults}, or \code{AllocHostLA} in
\code{GPU\_LAReference}) records which tier succeeded.

\paragraph{Transparent access.}
Because \code{cudaMallocHost} returns a \emph{pinned} host pointer
that is GPU-visible via unified addressing, CUDA kernels can read
and write the buffer with no code changes---only bandwidth
decreases (PCIe~4.0 $\approx$\,32\,GB/s versus device memory at
$\approx$\,500--1000\,GB/s).  For data that is streamed
sequentially (e.g.\ reference orbit look-ups during perturbation
iteration), the performance impact is manageable.

\paragraph{Matched cleanup.}
Destructors inspect the allocation-tier flag and call the
corresponding free function (\code{cudaFree} for device memory,
\code{cudaFreeHost} for pinned host memory), ensuring correct
resource release regardless of which tier was used.

The tiered allocation strategy introduced here addresses one aspect of memory
pressure.  The next chapter examines a complementary mechanism---disk-backed
growable vectors---that manages the large reference-orbit buffers whose size
can exceed available physical memory.
